%<*driver>
\def\nameofplainTeX{plain}
\ifx\fmtname\nameofplainTeX
\input l3docstrip.tex
\keepsilent
\askforoverwritefalse
\preamble

========================================================================
 Package tagdirtree - Accessible directory-tree listings.

 (C) 2026 Matthias Werner
 E-mail: matthias.werner@informatik.tu-chemnitz.de

 Released under the LaTeX Project Public License v1.3c or later
 See http://www.latex-project.org/lppl.txt
========================================================================

\endpreamble
\postamble

 This work is "maintained" (as per LPPL maintenance status) by
 Matthias Werner.
\endpostamble
\generate{\file{tagdirtree.sty}{\from{tagdirtree.dtx}{package}}}
\expandafter\endbatchfile
\fi

\iffalse
%</driver>
%<*package>
\NeedsTeXFormat{LaTeX2e}[2022-06-01]
\ProvidesExplPackage{tagdirtree}{2026-09-14}{v0.1.0}
  {Accessible directory-tree listings}
%</package>
%<*driver>
\fi
\documentclass[add-bib=false]{osgdoc}

\OnlyDescription
\usepackage{enumitem}
\usepackage{microtype}
\usepackage[
languages={de,en},
map={de=german,en=english},
targetlang={job=2},
load babel
]{langselect}
\newpackagename{\tagdirtreepkg}{tagdirtree}
\makeindex
\setcnltx{
  package  =  tagdirtree,
  version  = \UseVersionOf{tagdirtree.sty},
  date     = \UseDateOf{tagdirtree.sty},
  title    = {\ldeen{Das \tagdirtreepkg\ Paket}{The \tagdirtreepkg\ Package}},
  info     = {\ldeen{Verzeichnisbäume mit echter, nicht nur nachempfundener
      Barrierefreiheit}{Directory-tree listings with real, not merely
      simulated, accessibility}},
  authors  = {Matthias Werner[matthias.werner@informatik.tu-chemnitz.de]},
  abstract={
    \ldeen*{
      Anders als bestehende Verzeichnisbaum-Pakete (\pkg*{dirtree},
      \pkg*{dirtreex}) zeichnet \tagdirtreepkg\ den Baum nicht als
      eigenständige Grafik, die anschließend eine Alt-Text-Beschreibung
      braucht, sondern setzt ihn von Anfang an als echte, verschachtelte
      Liste: jeder Eintrag ist regulärer, lesbarer Text, die Verschachtelung
      selbst trägt die Eltern-Kind-Hierarchie, und pro Eintrag wird
      zusätzlich vermerkt, ob es sich um eine Datei, ein Verzeichnis oder
      ein Archiv handelt. Ist \cs{DocumentMetadata}\Marg{tagging=on}
      gesetzt, wird das über eine Alt-Beschreibung je Eintrag (Art, Name,
      Anmerkung) zusätzlich für unterstützende Technologie verfügbar
      gemacht; ohne Tagging bleibt es eine gewöhnliche verschachtelte
      Liste.
    }{
      Unlike existing directory-tree packages (\pkg*{dirtree},
      \pkg*{dirtreex}), \tagdirtreepkg\ does not draw the tree as a
      self-contained graphic that subsequently needs an alt-text
      description -- it sets it as a genuine nested list from the start:
      every entry is ordinary, readable text, the nesting itself carries
      the parent-child hierarchy, and each entry additionally records
      whether it is a file, a directory, or an archive. When
      \cs{DocumentMetadata}\Marg{tagging=on} is set, this is further
      exposed to assistive technology via a per-entry alt description
      (kind, name, note); without tagging it remains an ordinary nested
      list.
    }
  },
  url      =https://github.com/werner-matthias/osglecture,
  build-title
}
\begin{document}

\section{\ldeen{Zweck und Nutzung}{Purpose and usage}}

\ldeen{Das Paket hat keine Paketoptionen. Die Umgebung
\env{tagdirtree} rahmt einen Baum; \cs{dtdir}, \cs{dtfile} und
\cs{dtarchive} tragen jeweils Verzeichnisse, Dateien und Archive ein.
\cs{dtdir} nimmt als drittes Argument die Einträge des Verzeichnisses.}
{The package has no package options. The \env{tagdirtree} environment
frames a tree; \cs{dtdir}, \cs{dtfile}, and \cs{dtarchive} enter
directories, files, and archives respectively. \cs{dtdir} takes the
directory's entries as its third argument.}

\begin{verbatim}
\begin{tagdirtree}
  \dtdir{project}{root directory}{
    \dtdir{src}{source code}{
      \dtfile{main.py}{entry point}
      \dtfile{util.py}{}
    }
    \dtfile{README.md}{documentation}
    \dtarchive{legacy.zip}{old sources}
  }
\end{tagdirtree}
\end{verbatim}

\ldeen{Die zweite (bei \cs{dtdir} die zweite) Klammer ist die Anmerkung
zum Eintrag und darf leer sein. Sie erscheint im Text in Klammern und,
sofern gesetzt, als Teil der Alt-Beschreibung.}{The second brace
argument (also the second for \cs{dtdir}) is the entry's note and may
be empty. It appears in the text in parentheses and, when set, as part
of the alt description.}

\ldeen{Verzeichnisse werden fett, Archive kursiv, Dateien in
Grundschrift gesetzt; Einträge sind eine gewöhnliche (mit
\pkg*{enumitem} eng gesetzte) verschachtelte Liste, keine eigene
Zeichnung -- Optik lässt sich daher mit den üblichen
\pkg*{enumitem}-Mitteln auf der \env{tagdirtree}-Umgebung selbst
anpassen.}{Directories are set bold, archives italic, files in the
body font; entries form an ordinary (via \pkg*{enumitem} tightly set)
nested list, not a custom drawing -- appearance can therefore be
adjusted with the usual \pkg*{enumitem} means on the \env{tagdirtree}
environment itself.}

\subsection{\ldeen{Icons}{Icons}}

\ldeen{Jeder Eintrag setzt vor Name und Anmerkung ein kurzes,
sprachunabhängiges Icon -- \cs{tagdirtreefileicon},
\cs{tagdirtreedirectoryicon} bzw.\ \cs{tagdirtreearchiveicon}, je nach
Art. Die Vorgabe ist ein reiner Platzhalter (\code{[D]}, \code{[F]},
\code{[A]}) und zum Überschreiben gedacht; ist Tagging aktiv, wird das
Icon als rein dekoratives \code{artifact} markiert, da die Art des
Eintrags bereits vollständig im \code{alt}-Text steckt (siehe unten) --
ein Screenreader soll sie nicht zusätzlich vorlesen.}{Every entry sets
a short, language-independent icon before its name and note --
\cs{tagdirtreefileicon}, \cs{tagdirtreedirectoryicon}, or
\cs{tagdirtreearchiveicon}, depending on kind. The default is a plain
placeholder (\code{[D]}, \code{[F]}, \code{[A]}) and meant to be
overridden; while tagging is active, the icon is marked as a purely
decorative \code{artifact}, since the entry's kind already lives
entirely in its \code{alt} text (see below) -- a screen reader should
not read it out a second time.}

\begin{verbatim}
\RenewDocumentCommand \tagdirtreedirectoryicon { } { \faFolder }
\RenewDocumentCommand \tagdirtreefileicon { } { \faFile }
\RenewDocumentCommand \tagdirtreearchiveicon { } { \faFileArchive }
\end{verbatim}

\ldeen{Mit \pkg*{osgstyler} lassen sich Icons stattdessen in einer
Themefarbe statt einer festen Farbe setzen, über das öffentliche
\cs{OsgColorName}. (\pkg*{osgstyler}s \code{affix}-Decoration selbst
wurde geprüft und eignet sich \emph{nicht} als Anschlussstelle: Sie ist
an \pkg*{osgstyler-block}s eigene Anwendungsstellen gebunden, hat kein
öffentliches "`wende einen einzelnen Affix freistehend an"'-Kommando,
und ihr \code{symbol=}-Schlüssel hat ohnehin kein auflösendes Backend --
\code{content=} bliebe die einzig nutzbare Option, ohne echten
Mehrwert gegenüber direktem \cs{OsgColorName}.)}{With
\pkg*{osgstyler}, icons can instead be set in a theme color rather
than a fixed one, via the public \cs{OsgColorName}. (\pkg*{osgstyler}'s
own \code{affix} decoration was checked and does \emph{not} work as an
attachment point: it is bound to \pkg*{osgstyler-block}'s own
application sites, has no public "`apply a single affix
standalone"' command, and its \code{symbol=} key has no resolving
backend anyway -- \code{content=} would be the only usable option
left, with no real benefit over \cs{OsgColorName} directly.)}

\begin{verbatim}
\RenewDocumentCommand \tagdirtreedirectoryicon { }
  { \textcolor{\OsgColorName{accent}}{[D]} }
\end{verbatim}

\section{Implementation}
\MaybeStop{\ldeen*{Um die Implementationsbeschreibung zu erhalten, kommentieren Sie bitte \cs{OnlyDescription} in @1 aus.}{To view the implementation description, please uncomment \cs{OnlyDescription} in @1.}{\File{tagdirtree.dtx}}}

\AutoDocInput{tagdirtree.dtx}{package}
\end{document}

%</driver>

%<*package>
\RequirePackage{enumitem}
\ExplSyntaxOn

% \ldeen*{Anders als @1/@2, die den Baum als eigenständige Grafik
% zeichnen (roh bzw.\ über \pkg{tikz}) und daher -- ungetaggt -- als
% anonymes Artifact ohne jede Semantik im Tag-Baum landen, ist ein
% Verzeichnisbaum bereits von Natur aus eine verschachtelte Liste. @3
% baut ihn deshalb direkt als solche (\env{itemize} über \pkg*{enumitem})
% -- die Eltern-Kind-Hierarchie steckt in der Verschachtelung selbst,
% nicht in gezeichneten Verbindungslinien, und jeder Eintrag ist echter,
% lesbarer Text statt Teil eines Bildes. Bei aktivem Tagging bekommt
% jeder Eintrag zusätzlich eine Alt-Beschreibung (Art, Name, Anmerkung),
% da die für den Nutzer tatsächlich hörbare Semantik in der Praxis über
% diesen Kanal ankommt, nicht über eine eigene Strukturrolle (siehe
% unten bei @4).}{Unlike @1/@2, which draw the tree as a self-contained
% graphic (raw, resp.\ via \pkg{tikz}) and therefore -- once tagged --
% land as an anonymous artifact with no semantics at all in the tag
% tree, a directory tree already is a nested list by nature. @3
% therefore builds it directly as one (\env{itemize} via
% \pkg*{enumitem}) -- the parent-child hierarchy lives in the nesting
% itself, not in drawn connector lines, and every entry is real,
% readable text rather than part of a picture. While tagging is active,
% every entry additionally gets an alt description (kind, name, note),
% since that is the channel through which the semantics actually reach
% the user in practice, not a custom structure role (see @4
% below).}{\pkg*{dirtree}}{\pkg*{dirtreex}}{\tagdirtreepkg}{\cs{__tagdirtree_tag_begin:}}
\cs_if_exist:NT \NewStructureName
  {
    \NewStructureName{tagdirtree/file}
    \AssignStructureRole{tagdirtree/file}{Span}
    \NewStructureName{tagdirtree/directory}
    \AssignStructureRole{tagdirtree/directory}{Span}
    \NewStructureName{tagdirtree/archive}
    \AssignStructureRole{tagdirtree/archive}{Span}
  }

\tl_new:N \l__tagdirtree_tag_tl
\tl_new:N \l__tagdirtree_alt_tl

% \ldeen*{@1/@2 tragen die Rolle Span (nicht etwa LI): Der umgebende
% \env{itemize}-Eintrag ist bereits vom Kernel selbst korrekt als
% Listenpunkt getaggt (@3/@4), ein zweites, manuell geöffnetes @5 an
% derselben Stelle wäre ein ungültiges Eltern-Kind-Verhältnis
% (verifiziert: erzeugt \code{Relation is not allowed}-Fehler). @6
% direkt nach \cs{item} ist ebenfalls notwendig, nicht kosmetisch --
% ohne es öffnet die manuelle Tag-Struktur zu früh, bevor der Kernel
% die Absatzstruktur des Eintrags selbst begonnen hat, und landet unter
% dem falschen (noch nicht aktuellen) Elternknoten (verifiziert:
% dieselbe \code{Relation is not allowed}-Fehlerklasse, unabhängig von
% der gewählten Rolle).}{@1/@2 carry the role Span (not LI): the
% surrounding \env{itemize} entry is already correctly tagged as a list
% item by the kernel itself (@3/@4), a second, manually opened @5 at the
% same spot would be an invalid parent-child relation (verified: produces
% a \code{Relation is not allowed} error). @6 right after \cs{item} is
% likewise necessary, not cosmetic -- without it the manual tag structure
% opens too early, before the kernel has started the entry's own
% paragraph structure, and ends up under the wrong (not-yet-current)
% parent (verified: the same \code{Relation is not allowed} error class,
% regardless of the chosen role).}{\cs{tagstructbegin}}{\cs{tagmcbegin}}
% {\code{item}}{\code{itembody}}{LI}{\cs{leavevmode}}
\cs_new_protected:Npn \__tagdirtree_tag_begin:
  {
    \cs_if_exist:NT \tagstructbegin
      {
        \exp_args:Ne \tagstructbegin
          { tag=\l__tagdirtree_tag_tl, alt={\l__tagdirtree_alt_tl} }
        \exp_args:Ne \tagmcbegin { tag=\l__tagdirtree_tag_tl }
      }
  }
\cs_new_protected:Npn \__tagdirtree_tag_end:
  {
    \cs_if_exist:NT \tagstructend
      { \tagmcend \tagstructend }
  }

% \ldeen*{Sprachwahl mit derselben Rückfallkette wie andernorts im
% \pkg*{osglecture}-Bundle (@1, sonst @2, sonst @3) -- beide Integrationen
% sind rein optional und existenzgeprüft, das Paket funktioniert
% eigenständig auch ohne \pkg*{langselect} oder \pkg*{osglecture}, dann
% stets auf @3.}{Language selection with the same fallback chain used
% elsewhere in the \pkg*{osglecture} bundle (@1, else @2, else @3) --
% both integrations are purely optional and existence-checked, the
% package works standalone without \pkg*{langselect} or \pkg*{osglecture}
% too, always falling back to @3 then.}
% {\cs{olsTargetLanguage}}{\cs{OsgLectureLanguage}}{\code{en}}
\cs_new:Npn \__tagdirtree_kind_word:n #1
  {
    \str_if_eq:eeTF
      {
        \cs_if_exist:NTF \olsTargetLanguage
          { \olsTargetLanguage }
          {
            \cs_if_exist:NTF \OsgLectureLanguage
              { \OsgLectureLanguage }
              { en }
          }
      }
      { de }
      {
        \str_case:nn {#1}
          {
            { file }      { Datei }
            { directory } { Verzeichnis }
            { archive }   { Archiv }
          }
      }
      {
        \str_case:nn {#1}
          {
            { file }      { File }
            { directory } { Directory }
            { archive }   { Archive }
          }
      }
  }

% \ldeen*{@1--@3 sind bewusst gewöhnliche, überschreibbare Befehle statt
% eines Schlüssel-Wert-Optionssystems -- ihr Standardinhalt ist ein
% reiner Platzhalter (\code{[D]} usw.), gedacht zum direkten
% Überschreiben (\cs{renewcommand} oder \cs{RenewDocumentCommand}),
% etwa mit \pkg*{osgstyler}-Decorations (\code{content=}, nicht
% \code{symbol=} -- Letzteres hat noch keinen auflösenden Backend). Das
% Icon selbst wird als \code{artifact} markiert, da die für das
% Verständnis nötige Art-Angabe bereits vollständig im \code{alt}-Text
% des Eintrags steckt (siehe oben); ohne diese Markierung würde ein
% Screenreader Icon und Alt-Text redundant nacheinander
% ankündigen.}{@1--@3 are deliberately ordinary, overridable commands
% rather than a key-value option system -- their default content is a
% plain placeholder (\code{[D]} etc.), meant to be overridden directly
% (\cs{renewcommand} or \cs{RenewDocumentCommand}), e.g.\ with
% \pkg*{osgstyler} decorations (\code{content=}, not \code{symbol=} --
% the latter has no resolving backend yet). The icon itself is marked
% as \code{artifact}, since the kind information needed for
% understanding already lives entirely in the entry's \code{alt} text
% (see above); without this marking a screen reader would announce icon
% and alt text redundantly, one after the other.}
% {\cs{tagdirtreefileicon}}{\cs{tagdirtreedirectoryicon}}
% {\cs{tagdirtreearchiveicon}}
\NewDocumentCommand \tagdirtreefileicon { } { [F] }
\NewDocumentCommand \tagdirtreedirectoryicon { } { [D] }
\NewDocumentCommand \tagdirtreearchiveicon { } { [A] }

% \ldeen*{@1 zählt die umgebenden Verzeichnisse mit (0 auf oberster
% Ebene); es ist ein reiner Zähler, keine Positionsinformation
% innerhalb der Geschwisterliste -- @2 (Standard: leer) kann daraus
% zwar eine tiefenabhängige, aber keine "`letztes Element"'-bewusste
% Optik bauen (kein \code{├} gegenüber \code{└}), da dafür bekannt sein
% müsste, wie viele weitere Geschwister noch folgen, was ohne
% Vorausschau über das jeweilige \cs{dtdir}-Textargument hinweg nicht
% zur Verfügung steht. @3 zeigt, wie sich damit trotzdem eine
% klassische Verbindungslinien-Optik bauen lässt, samt Kopplung an
% \pkg*{osgstyler}s Themefarbe -- bewusst nur als Beispiel, nicht als
% zusätzliche Paketabhängigkeit.}{@1 counts the surrounding directories
% (0 at the top level); it is a plain counter, not positional
% information within the sibling list -- @2 (default: empty) can build
% a depth-aware look from it, but not a "`last item"'-aware one (no
% \code{├} versus \code{└}), since that would require knowing how many
% further siblings still follow, which is not available without
% look-ahead across the respective \cs{dtdir} text argument. @3 shows
% how a classic connector-line look can still be built from this,
% including coupling to \pkg*{osgstyler}'s theme color -- deliberately
% only as an example, not as an added package dependency.}
% {\cs{tagdirtreedepth}}{\cs{tagdirtreeconnector}}
% {\File{tagdirtree-lines-skin.tex}}
\int_new:N \g__tagdirtree_depth_int
\NewExpandableDocumentCommand \tagdirtreedepth { }
  { \int_use:N \g__tagdirtree_depth_int }
\NewDocumentCommand \tagdirtreeconnector { } { }

\cs_new:Npn \__tagdirtree_icon:n #1
  {
    \str_case:nn {#1}
      {
        { file }      { \tagdirtreefileicon }
        { directory } { \tagdirtreedirectoryicon }
        { archive }   { \tagdirtreearchiveicon }
      }
  }

\cs_new_protected:Npn \__tagdirtree_render:nnn #1#2#3
  {
    \cs_if_exist:NT \tagmcbegin { \tagmcbegin{artifact} }
    \tagdirtreeconnector
    \__tagdirtree_icon:n {#1}
    \cs_if_exist:NT \tagmcend { \tagmcend }
    ~
    \str_case:nnF {#1}
      {
        { directory } { \textbf{#2} }
        { archive }   { \textit{#2} }
      }
      { #2 }
    \tl_if_empty:nF {#3} { \space (#3) }
  }

\cs_new_protected:Npn \__tagdirtree_entry:nnnn #1#2#3#4
  {
    \item
    \leavevmode
    \tl_set:Ne \l__tagdirtree_tag_tl { \UseStructureName{#2} }
    \tl_if_empty:nTF {#4}
      { \tl_set:Ne \l__tagdirtree_alt_tl { \__tagdirtree_kind_word:n{#1}:~#3 } }
      {
        \tl_set:Ne \l__tagdirtree_alt_tl
          { \__tagdirtree_kind_word:n{#1}:~#3~--~#4 }
      }
    \__tagdirtree_tag_begin:
    \__tagdirtree_render:nnn {#1}{#3}{#4}
    \__tagdirtree_tag_end:
  }

\NewDocumentCommand \dtfile { m m }
  { \__tagdirtree_entry:nnnn {file} {tagdirtree/file} {#1}{#2} }
\NewDocumentCommand \dtarchive { m m }
  { \__tagdirtree_entry:nnnn {archive} {tagdirtree/archive} {#1}{#2} }
\NewDocumentCommand \dtdir { m m +m }
  {
    \__tagdirtree_entry:nnnn {directory} {tagdirtree/directory} {#1}{#2}
    \int_gincr:N \g__tagdirtree_depth_int
    \begin{itemize}
      #3
    \end{itemize}
    \int_gdecr:N \g__tagdirtree_depth_int
  }

\NewDocumentEnvironment{tagdirtree}{}
  { \begin{itemize}[nosep] }
  { \end{itemize} }

\ExplSyntaxOff
%</package>
